\chapter*{Eidesstattliche Erklärung}
\label{ch:erklaerung}

\begin{flushleft}
Ich erkläre an Eides statt, dass
\end{flushleft}

\begin{flushleft}
- ich die vorliegende Arbeit selbständig und ohne fremde Hilfe verfasst, keine anderen als die angegebenen Quellen und Hilfsmittel benutzt habe.
\end{flushleft}

\begin{flushleft}
- ich mich bei der Erstellung der Arbeit an die Standards guter wissenschaftlicher Praxis gemäß dem Leitfaden zum Wissenschaftlichen Arbeiten der FH St. Pölten gehalten habe.
\end{flushleft}

\begin{flushleft}
- ich die vorliegende Arbeit an keiner Hochschule zur Beurteilung oder in irgendeiner Form als Prüfungsarbeit vorgelegt oder veröffentlicht habe. 
\end{flushleft}

Über den Einsatz von Hilfsmitteln der generativen Künstlichen Intelligenz wie Chatbots, Bildgeneratoren, Programmieranwendungen, Paraphrasier- oder Übersetzungstools erkläre ich, dass

% Durch das Kommando [\done] nach dem \item können die einzelnen Checkboxen abgehakt werden.

\begin{todolist}
    \item[\done] im Zuge dieser Arbeit kein Hilfsmittel der generativen Künstlichen Intelligenz zum Einsatz gekommen ist. 
    \item ich Hilfsmittel der generativen Künstlichen Intelligenz verwendet habe, um die Arbeit Korrektur zu lesen.
    \item ich Hilfsmittel der generativen Künstlichen Intelligenz verwendet habe, um Teile des Inhalts der Arbeit zu erstellen. Ich versichere, dass ich jeden generierten Inhalt mit der Originalquelle zitiert habe. Das genutzte Hilfsmittel der generativen Künstlichen Intelligenz ist an entsprechenden Stellen ausgewiesen.
\end{todolist}

Durch den Leitfaden zum Wissenschaftlichen Arbeiten der FH St. Pölten bin ich mir über die Konsequenzen einer wahrheitswidrigen Erklärung bewusst.

\vspace{1.5cm}


